%% ma: L11-B16-0098
%% lop: 11
%% chuong: 5
%% bai: 16
%% dang: 11.16.4
%% loai: TL
%% muc: VDC
%% dap_an: "-"
%% nguon: "Ảnh giáo viên gửi 25/11/2026: Giới hạn hàm số dạng 0/0 (3 bài, có ĐS), Bài 3 câu 6"
%% kien_thuc: "Lớp 11 Bài 16 (giới hạn của hàm số); hằng đẳng thức x^k-1=(x-1)(x^(k-1)+...+1); tổng 1+2+...+n (Lớp 11 Bài 6 cấp số cộng)"
%% trang_thai: chinh-thuc
%% ghi_chu: "Đề gốc không ghi điều kiện của m, n, đã thêm m, n là số nguyên dương."
\begin{ex}
Tính giới hạn $A=\lim\limits_{x\to1}\left(\dfrac{m}{1-x^m}-\dfrac{n}{1-x^n}\right)\ (m,n\text{ là số nguyên dương})$.
\loigiai{Ta có \[\dfrac m{1-x^m}-\dfrac1{1-x}=\dfrac{(1-x)+(1-x^2)+\dots+(1-x^{m-1})}{1-x^m}\] \[=\dfrac{1+(1+x)+\dots+(1+x+\dots+x^{m-2})}{1+x+\dots+x^{m-1}},\] dần về $\dfrac{1+2+\dots+(m-1)}{m}=\dfrac{m-1}2$. Tương tự với $n$. Vậy $A=\dfrac{m-1}2-\dfrac{n-1}2=\dfrac{m-n}2$.}
\end{ex}
